% ======================================================================
% Capítulo: Estado del arte y tecnologías
% Fundamentos de LLM y RAG, chunking, bases vectoriales, evaluación y tecnologías. Extensión orientativa: 7--9 páginas.
% ======================================================================
\chapter{Estado del arte y tecnologías}
\label{cap:estado-arte}

\section{Modelos de lenguaje y arquitectura Transformer}
\label{sec:llm}
% TODO: Capacidades, ventana de contexto y coste \cite{vaswani2017}.

\section{Alucinaciones en LLM}
\label{sec:alucinaciones}
% TODO: Taxonomía y por qué son críticas en el dominio jurídico \cite{huang2023}.

\section{Retrieval-Augmented Generation}
\label{sec:rag}
% TODO: Esquema general y variantes: \emph{naive}, \emph{advanced}, jerárquico \emph{small-to-big} \cite{lewis2020,gao2023}.

\section{Estrategias de fragmentación (\emph{chunking})}
\label{sec:chunking}
% TODO: Fixed-size con solapamiento, recursive character splitting, semantic chunking y padre/hijo.

\section{Embeddings y bases de datos vectoriales}
\label{sec:vectoriales}
% TODO: Similitud coseno, ANN (HNSW, IVFFlat), pgvector frente a Qdrant \cite{wang2024,malkov2020}.

\section{Recuperación híbrida}
\label{sec:hibrida}
% TODO: Búsqueda densa, búsqueda léxica (\texttt{tsvector}) y fusión RRF \cite{cormack2009}.

\section{Evaluación de sistemas RAG}
\label{sec:evaluacion-rag}
% TODO: RAGAS \cite{es2023ragas}, TruLens, métricas de fidelidad y contexto; evaluación por preferencia humana \cite{chiang2024}.

\section{Orquestación de IA}
\label{sec:orquestacion}
% TODO: Semantic Kernel \cite{semantickernel} frente a LangChain y LlamaIndex.

\section{Trabajos relacionados}
\label{sec:relacionados}
% TODO: Asistentes de IA en la administración pública y herramientas jurídicas; qué aporta este TFG.

\section{Tecnologías empleadas}
\label{sec:tecnologias}
% TODO: Tabla con cada tecnología y el motivo de su elección (ver tabla de ejemplo abajo).

\begin{table}[ht]
  \centering
  \caption{Tecnologías empleadas y motivo de su elección.}
  \label{tab:tecnologias}
  \begin{tabularx}{\linewidth}{lX}
    \toprule
    Tecnología & Motivo de la elección \\
    \midrule
    .NET 10 y Aspire & % TODO
    \\
    Go & % TODO
    \\
    Semantic Kernel & % TODO
    \\
    PostgreSQL y Qdrant & % TODO
    \\
    RabbitMQ & % TODO
    \\
    Angular & % TODO
    \\
    \bottomrule
  \end{tabularx}
\end{table}
